
\paragraph{A Simple Baseline: Using Extracted Mentions as Captions.}
%\kenneth{MAYBE not put it here...}
%As suggested by prior literature that images usually illustrate objects in the surrounding documents~\cite{6205758} and 
Motivated by our information overlap study (\Cref{sec:motivating-analysis}),
we created the \textbf{Reuse} baselines. 
These baselines simply repurpose portions of the input text as the prediction.
%we designed the \textbf{Reuse} baselines that simply output parts of the input text (\eg, texts surrounding Mention) as the prediction. 
% For example, Reuse$_M$ (Reuse with Mention) outputs the first mention as the caption;
% Reuse$_{W[0, 1]}$ (Reuse with Window[0, 1]) outputs the first mention and the following 1 sentence as the caption.
% \kenneth{CY TODO: Check if it makes sense.}
% and Reuse$_{P}$ (Reuse with Paragraph) outputs the whole paragraph as the summary.
%Similar to prior findings that relevant events are more likely to occur closely~\cite{huang-huang-2021-semantic}, 
%we expected Reuse$_M$ to perform the best among the Reuse baselines and its performance would drop as the window size increases.

% \kenneth{We need to cite some more papers to justify this baseline. (1) Semantic Frame Forcast paper and (2) maybe this \cite{6205758}}
% \cy{Done}

% \paragraph{Baseline: Vision-to-Language Generation.}
\paragraph{Vision-to-Language Baselines.}
The vision-to-language generation treated this task as an image-captioning task 
that took the scientific figure image as input and generated a text to describe it.
%For this baseline, 
We compared two vision-to-language models as baselines.
First, we built a sequence-to-sequence model by combining BEiT~\cite{bao2022beit} and GPT-2~\cite{radford2019language}.
%The second model we chose was the TrOCR~\cite{li2021trocr} model, a transformer-based sequence-to-sequence model pre-trained for the OCR task.
%Compared to image encoders trained on real photos (images containing visual objects) 
%such as ViT~\cite{dosovitskiy2020vit} and BEiT~\cite{bao2022beit}, 
%OCR models trained with images of printed and handwritten documents are closer to the scientific paper domain, 
%where figures are mostly lines and texts.
We also selected the TrOCR~\cite{li2021trocr} model, a transformer-based sequence-to-sequence model pre-trained for OCR tasks.
Compared to image encoders like ViT~\cite{dosovitskiy2020vit} and BEiT~\cite{bao2022beit}, which were trained on photos, OCR models trained on printed and handwritten documents align more closely with the scientific paper domain.
All figures from \oldDataset (106,391 training samples) were used for training since no mentions were required.
%Note that 
%All available figures from \oldDataset (106,391 training samples) were used for training, as no mentions were needed.
%\kenneth{Is is before or after resplit?}\cy{After the resplit.}
% HuggingFace's implementation of the TrOCR model\footnote{\texttt{microsoft/trocr-large-printed} was used.} was used for fine-tuning.
% The model was trained for 100 epochs, and the one with the highest ROUGE-2 score was kept.

% \paragraph{Baseline: [X-Axis] and [Y-Axis]}

% \kenneth{TODO ED}